\documentclass[12pt,a4paper]{article}

%=========================================================
%                    ENCODAGE ET POLICES
%=========================================================
\usepackage[utf8]{inputenc}
\usepackage[T1]{fontenc}
\usepackage[french]{babel}
\usepackage{lmodern}

%=========================================================
%                      MATHÉMATIQUES
%=========================================================
\usepackage{amsmath,amssymb,amsfonts,mathrsfs}
\usepackage{tkz-tab}
\everymath{\displaystyle}

%=========================================================
%                      MISE EN PAGE
%=========================================================
\usepackage[left=1.5cm,right=1.5cm,top=1.5cm,bottom=1.8cm]{geometry}
\usepackage{multicol}
\usepackage{float}
\usepackage{array,multirow}
\usepackage{enumitem}

%=========================================================
%                 GRAPHISMES ET COULEURS
%=========================================================
\usepackage{tikz}
\usetikzlibrary{arrows.meta,calc}
\usepackage{pgfplots}
\pgfplotsset{compat=1.15}

\usepackage[most]{tcolorbox}
\tcbuselibrary{skins,hooks}

%=========================================================
%                  LIENS ET ARRIÈRE-PLAN
%=========================================================
\usepackage[colorlinks=true,linkcolor=blue,urlcolor=blue,citecolor=blue]{hyperref}
\usepackage{eso-pic}

%=========================================================
%                    COULEURS ET STYLES
%=========================================================
\definecolor{myred}{RGB}{180, 40, 40}
\definecolor{myblue}{RGB}{20, 50, 120}
\definecolor{mygreen}{RGB}{30, 120, 40}

%=========================================================
%                        FILIGRANE
%=========================================================
\AddToShipoutPicture{
    \AtTextCenter{
        \makebox[0pt]{
            \rotatebox{45}{
                \textcolor[gray]{0.92}{
                    \fontsize{5cm}{5cm}\selectfont PGB
                }
            }
        }
    }
}

\begin{document}

% En-tête
\begin{center}
    \Large\textbf{\underline{\textcolor{myred}{CHAPITRE 7 : FONCTION EXPONENTIELLE NÉPÉRIENNE}}}\\[0.2cm]
    \normalsize\textbf{Prof : M. BA} \hfill \textbf{Classe : Terminale L}\\[0.1cm]
    \textbf{Durée : 6 heures} \hfill \textbf{Année scolaire : 2025 -- 2026}
\end{center}
\hrule
\vspace{0.3cm}

%=========================================================
%    I. Généralités
%=========================================================

\section*{\color{myred}I. Généralités}

\paragraph{\underline{Activité 1}}
À l'aide de la calculatrice, compléter le tableau suivant (donner les valeurs arrondies à $10^{-2}$ près) :

\begin{center}
\renewcommand{\arraystretch}{1.5}
\begin{tabular}{|c|c|c|c|c|c|c|c|}
\hline
$x$ & -2 & -1 & 0 & 0,5 & 1 & 2 & 3 \\
\hline
$e^x$ & \hspace*{0.8cm} & \hspace*{0.8cm} & \hspace*{0.8cm} & \hspace*{0.8cm} & \hspace*{0.8cm} & \hspace*{0.8cm} & \hspace*{0.8cm} \\
\hline
\end{tabular}
\end{center}

\vspace{0.3cm}

\subsection*{\color{myred}1. Définition et notation}
\begin{tcolorbox}[colback=blue!5!white,colframe=myblue,title=\textbf{Définition}]
La fonction logarithme népérien $\ln$ réalisant une bijection de $]0\,;\,+\infty[$ sur $\mathbb{R}$, elle admet une fonction réciproque définie sur $\mathbb{R}$ à valeurs dans $]0\,;\,+\infty[$. 

Cette fonction réciproque est appelée \textbf{fonction exponentielle népérienne}, notée $\exp$ ou $x \mapsto e^x$.
\[
\text{Pour tout } x \in \mathbb{R} \text{ et pour tout } y \in ]0\,;\,+\infty[ : \quad e^x = y \iff x = \ln(y)
\]
\end{tcolorbox}

\vspace{0.3cm}

\subsection*{\color{myred}2. Remarques et Relations fondamentales}

\paragraph{\underline{Remarque 1}}
\begin{itemize}
    \item La fonction $e^x$ est définie sur $\mathbb{R}$ tout entier ($D_{\exp} = \mathbb{R}$).
    \item Pour tout réel $x$, $\mathbf{e^x > 0}$ (la fonction exponentielle est **strictement positive**).
    \item Valeurs particulières : $e^0 = 1$ et $e^1 = e \approx 2{,}718$.
\end{itemize}

\vspace{0.3cm}

\paragraph{\underline{Remarque 2 (Simplifications de réciprocité)}}
Pour tout réel $x$ et pour tout réel $y > 0$ :
\begin{itemize}
    \item $\ln(e^x) = x \quad (\text{pour tout } x \in \mathbb{R})$
    \item $e^{\ln(y)} = y \quad (\text{pour tout } y > 0)$
\end{itemize}

\subsection*{\color{myred}3. Propriétés algébriques}

\begin{tcolorbox}[colback=red!5!white,colframe=myred,title=\textbf{Propriété fondamentale (Produit)}]
Pour tous réels $a$ et $b$ :
\[
e^{a + b} = e^a \times e^b
\]
\end{tcolorbox}

\paragraph{Conséquences des propriétés :}
Pour tous réels $a$ et $b$, et pour tout $n \in \mathbb{Z}$ :
\begin{enumerate}[label=\alph*)]
    \item \textbf{Inverse :} $e^{-a} = \dfrac{1}{e^a}$
    \item \textbf{Quotient :} $e^{a - b} = \dfrac{e^a}{e^b}$
    \item \textbf{Puissance :} $(e^a)^n = e^{n \cdot a}$
    \item \textbf{Racine carrée :} $e^{\frac{1}{2}a} = \sqrt{e^a}$
\end{enumerate}

\vspace{0.3cm}

\subsection*{\color{myred}Exemple :}
Simplifier l'expression $A = \dfrac{e^5 \times e^{-2}}{e^1} + (e^2)^3$.

\subsection*{\color{myred}Solution :}
$A = \dfrac{e^{5 - 2}}{e^1} + e^{2 \times 3} = \dfrac{e^3}{e^1} + e^6 = e^{3 - 1} + e^6 = e^2 + e^6$

\subsection*{\color{myred}Application :}
\begin{enumerate}
    \item Écrire $A = \dfrac{e^3 \times e^{-5}}{e^2}$ sous la forme $e^k$ où $k \in \mathbb{Z}$.
    \item Simplifier $B = (e^x + e^{-x})^2 - (e^x - e^{-x})^2$.
    \item Simplifier $C = \dfrac{e^{2x} - 1}{e^x + 1}$.
\end{enumerate}

\subsection*{\color{myred}Correction :}
\begin{enumerate}
\item Écrivons\\ $
\begin{aligned}
A &= \dfrac{e^3 \times e^{-5}}{e^2} \\[3mm]
  &= \dfrac{e^{3 - 5}}{e^2} = \dfrac{e^{-2}}{e^2} = e^{-2 - 2} = \mathbf{e^{-4}}
\end{aligned}
$
    \item Simplifions\\ $
\begin{aligned}
B &= (e^x + e^{-x})^2 - (e^x - e^{-x})^2\\[3mm]
  &= \left((e^x)^2 + 2e^x e^{-x} + (e^{-x})^2\right) - \left((e^x)^2 - 2e^x e^{-x} + (e^{-x})^2\right)\\[3mm]
  &= (e^{2x} + 2e^0 + e^{-2x}) - (e^{2x} - 2e^0 + e^{-2x})\\[3mm]
  &= e^{2x} + 2 + e^{-2x} - e^{2x} + 2 - e^{-2x} = \mathbf{4}
\end{aligned}
$

\item Simplifions\\ $
\begin{aligned}
C &= \dfrac{e^{2x} - 1}{e^x + 1} = \dfrac{(e^x)^2 - 1^2}{e^x + 1}\\[4mm]
  &= \dfrac{(e^x - 1)(e^x + 1)}{e^x + 1} = \mathbf{e^x - 1}
\end{aligned}
$
\end{enumerate}

%=========================================================
%    II. ÉQUATIONS ET INÉQUATIONS FAISANT INTERVENIR EXP
%=========================================================

\section*{\color{myred}II. Équations-Inéquations}
Grâce à la stricte croissance de la fonction exponentielle sur $\mathbb{R}$, pour tous réels $a$ et $b$ :
\begin{itemize}
    \item $e^a = e^b \iff a = b$
    \item $e^a < e^b \iff a < b$
    \item $e^x = k \iff x = \ln(k) \quad (\text{avec } k > 0)$
    \item $e^x \le 0 \implies \text{Impossible (pas de solution)}$
\end{itemize}

\subsection*{\color{myred}1. Équations}

\paragraph{Exemple 1 : Équation du type $e^{A(x)} = e^{B(x)}$}\mbox{}\\[2mm]
Résoudre dans $\mathbb{R}$ l'équation $(E_1) : e^{2x - 3} = e^{x + 1}$.

\begin{itemize}
    \item \textbf{Domaine de validité :}\\ $D_E = \mathbb{R}$.
    \item \textbf{Résolution :}\\
    Pour tout $x \in \mathbb{R}$ :
    \[
    e^{2x - 3} = e^{x + 1} \iff 2x - 3 = x + 1 \iff x = 4
    \]
    \item \textbf{Conclusion :}\\ $\mathbf{S = \{4\}}$.
\end{itemize}

\paragraph{Exemple 2 : Équation se ramenant au second degré}\mbox{}\\[2mm]
Résoudre dans $\mathbb{R}$ l'équation $(E_2) : e^{2x} - 3e^x + 2 = 0$.

\begin{itemize}
    \item \textbf{Domaine de validité :}\\ $D_E = \mathbb{R}$.
    \item \textbf{Changement de variable :}\\ Comme $e^{2x} = (e^x)^2$, posons $X = e^x$ avec $\mathbf{X > 0}$. L'équation devient :
    \[
    X^2 - 3X + 2 = 0
    \]
    $\Delta = (-3)^2 - 4(1)(2) = 1 > 0 \implies X_1 = 1 \quad \text{et} \quad X_2 = 2$.
    \item \textbf{Retour à $x$ :}
    \begin{itemize}
        \item $e^x = 1 \iff x = \ln(1) = 0$
        \item $e^x = 2 \iff x = \ln(2)$
    \end{itemize}
    \item \textbf{Conclusion :} $\mathbf{S = \{0 \,;\, \ln(2)\}}$.
\end{itemize}

\subsection*{\color{myred}2. Inéquations}

\paragraph{Exemple 3 :}\mbox{}\\[2mm]
Résoudre dans $\mathbb{R}$ l'inéquation $(I_1) : e^{3x - 1} \le 1$.

\begin{itemize}
    \item \textbf{Domaine de validité :}\mbox{}\\[2mm] $D_I = \mathbb{R}$.
    \item \textbf{Résolution :}\mbox{}\\[2mm]
    Comme $1 = e^0$, l'inéquation s'écrit :
    \[
    e^{3x - 1} \le e^0 \iff 3x - 1 \le 0 \iff 3x \le 1 \iff x \le \frac{1}{3}
    \]
    \item \textbf{Conclusion :}\mbox{}\\[2mm] $\mathbf{S = \Big]-\infty \,;\, \frac{1}{3}\Big]}$.
\end{itemize}

\vspace{0.3cm}
\hrule
\vspace{0.3cm}

\subsection*{\color{myred}Exercice d'application : Équations et inéquations}

\begin{tcolorbox}[colback=yellow!5!white,colframe=myred,title=\textbf{Exercice d'application}]
Résoudre dans $\mathbb{R}$ les équations et inéquations suivantes :
\begin{enumerate}
    \item $e^x = 1$
    \item $e^x < 2$
    \item $e^{x + 3} = 5$
\end{enumerate}
\end{tcolorbox}

\vspace{0.3cm}

\begin{tcolorbox}[colback=gray!5!white,colframe=myblue,title=\textbf{Correction détaillée}]

\paragraph{1. Résolution de $e^x = 1$ :}\leavevmode\\[2mm]
\begin{itemize}
    \item \textbf{Domaine de définition :} $D = \mathbb{R}$.
    \item \textbf{Résolution :}
    \[
    e^x = 1 \iff e^x = e^{\color{red}0} \iff x = \mathbf{0}
    \]
    \item \textbf{Conclusion :} L'ensemble des solutions est $\mathbf{S_1 = \{0\}}$.
\end{itemize}

\vspace{0.4cm}

\paragraph{2. Résolution de $e^x < 2$ :}\leavevmode\\[2mm]
\begin{itemize}
    \item \textbf{Domaine de définition :} $D = \mathbb{R}$.
    \item \textbf{Résolution :}
    
    La fonction exponentielle étant strictement croissante sur $\mathbb{R}$ :
    \[
    e^x < 2 \iff \ln(e^x) < \ln(2) \iff x < \mathbf{\ln(2)}
    \]
    \item \textbf{Conclusion :} L'ensemble des solutions est $\mathbf{S_2 = ]-\infty \,;\, \ln(2)[}$.
\end{itemize}

\vspace{0.4cm}

\paragraph{3. Résolution de $e^{x + 3} = 5$ :}\leavevmode\\[2mm]
\begin{itemize}
    \item \textbf{Domaine de définition :} $D = \mathbb{R}$.
    \item \textbf{Résolution :}
    \[
    \begin{aligned}
    e^{x + 3} = 5 &\iff x + 3 = \ln(5)\\[3mm]
                  &\iff x = \mathbf{\ln(5) - 3}
    \end{aligned}
    \]
    \item \textbf{Conclusion :} L'ensemble des solutions est $\mathbf{S_3 = \{\ln(5) - 3\}}$.
\end{itemize}

\end{tcolorbox}

\begin{tcolorbox}[colback=yellow!5!white,colframe=myred,title=\textbf{Devoir à la maison : Équations et inéquations exponentielles}]
Résoudre dans $\mathbb{R}$ les équations et inéquations suivantes :
\begin{enumerate}
    \item $e^{2x - 1} = 1$
    \item $2e^x - 3 \le 0$
    \item $e^{x - 4} = e^{2x - 3}$
    \item $e^{2x} + e^x - 2 \ge 0$
    \item $3(e^x)^2 - 2e^x - 1 = 0$
    \item $\dfrac{e^{2x - 1}}{e^{1 + x}} < e^3$
\end{enumerate}
\end{tcolorbox}
\end{document}